\documentclass[10pt,a4paper]{article}
\usepackage[utf8]{inputenc}
\usepackage[a4paper,left=3cm,right=2cm]{geometry}
\usepackage{amsmath}
\usepackage[thmmarks, amsmath, thref]{ntheorem}
\usepackage{amsfonts}
\usepackage{amssymb}
\usepackage{fancyhdr}
\pagestyle{fancy}
\usepackage{anysize}
\usepackage{graphicx}
\usepackage{float}
\usepackage{hyperref}
\usepackage{multirow}
\usepackage{enumitem}
\newtheorem{theorem}{Theorem}
\theoremsymbol{\ensuremath{\blacksquare}}
\newtheorem*{solution}{Solución:}
\usepackage{listings}
\usepackage{xcolor}

% Margenes y formalidades

\textwidth 6.25in % Margen Derecho y ancho del texto
\textheight 8.5in
\oddsidemargin 0in % Margen Izquierdo
\headheight 0in % Largo del texto
\leftmargin 1in
\parindent 0em % Salto entre la indentación (Sangría)https://www.overleaf.com/4519127669sdpptkhvsbfn
\topmargin -0.5in % Margen Superior
\parskip 0.5ex % Salto entre parrafos
\baselineskip 1pt

%\lhead{\includegraphics[scale=0.2]{logo.png}} % texto izquierda de la cabecera
%\chead{}                                      % texto centro de la cabecera
%\rhead{\includegraphics[scale=0.2]{FIG/logo-mba-2016.png}} % número de página a la derecha

\renewcommand{\headrulewidth}{0.4pt} % grosor de la línea de la cabecera
\renewcommand{\footrulewidth}{0.4pt} % grosor de la línea del pie

\begin{document}



\pagebreak
\vspace*{4mm} 

\begin{center}
\begin{huge}
\textbf{\\Ejercicios Econometría 4}
\end{Large}\\
%{\large 2023-2}\\
\end{center}

\begin{enumerate}

\item En una prueba estandarizada para postular a un cargo en especifico, se pudo encontrar, a través de una encuesta, las horas de estudio para esta prueba y la motivación por el cargo. Las notas, las horas de estudio y la motivación de 5 individuos están en la siguiente tabla:
  \[
    \begin{array}{|c|c|c|c|}
    \hline
      i & Notas (Y_i) & Horas (X_{1i}) & Motivación (X_{2i})\\\hline
      1 & 30 & 4  & 4 \\
      2 & 34 & 4  & 5 \\
      3 & 40 & 6  & 7 \\
      4 & 46 & 8  & 9 \\
      5 & 52 & 10 & 9\\
      \hline
    \end{array}
  \]

  \begin{enumerate}
    \item Estime el modelo \(Y_i=\beta_0+\beta_1X_{1i}+u_i\) e informe
      \(\hat{\beta}_0,\hat{\beta}_1, R^{2}, R^{2}_{\text{aj}}\) y
      \(\operatorname{se}(\hat{\beta}_1)\).
\\
\textbf{Solución}
\\
\[
\bar Y=40.4,\qquad
\bar X_1=6.4
\]

\begin{center}
\begin{tabular}{|c|r|r|r|r|r|r|r|}
\toprule
\hline
$i$ & $X_{1i}$ & $X_{1i}-\bar X_1$ & $Y_i$ & $Y_i-\bar Y$ &
$(X_{1i}-\bar X_1)^2$ & $(Y_i-\bar Y)^2$ & $(X_{1i}-\bar X_1)(Y_i-\bar Y)$\\
\midrule
\hline
1 & 4  & $-2.4$ & 30 & $-10.4$ &  5.76 & 108.16 & 24.96\\
2 & 4  & $-2.4$ & 34 & $-6.4$  &  5.76 &  40.96 & 15.36\\
3 & 6  & $-0.4$ & 40 & $-0.4$  &  0.16 &   0.16 &  0.16\\
4 & 8  & $ 1.6$ & 46 & $ 5.6$  &  2.56 &  31.36 &  8.96\\
5 & 10 & $ 3.6$ & 52 & $11.6$  & 12.96 & 134.56 & 41.76\\
\midrule
\hline
$\sum$&   & 0 &   & 0 & 27.20 & 315.20 & 91.20\\
\bottomrule
\end{tabular}
\end{center}


\[
\hat\beta_1=\frac{\sum (X_{1i}-\bar X_1)(Y_i-\bar Y)}{\sum (X_{1i}-\bar X_1)^2}=\frac{91.2}{27.2}=3.3529,
\qquad
\hat\beta_0=\bar Y-\hat\beta_1\bar X_1
          =40.4-3.3529\cdot6.4=18.9412.
\]
De la tabla ya sabemos que $SST = 315.2$, entonces encontrando $SSR$, tenemos:
\begin{center}
\begin{tabular}{|c|r|r|r|r|r|r|}
\toprule
\hline
$i$ & $X_{1i}$ & $Y_i$ & 
 $(Y_i-\bar Y)^2$ & $\hat{Y_i}$ & $\hat{u_i}=Y_i - \hat{Y_i}$ & $\hat{u_i}^2$\\
\midrule
\hline
1 & 4  &  30   & 108.16 & 32.35 & -2.353 & 5.536\\
2 & 4  &  34   &  40.96 & 32.35 & 1.647  & 2.713\\
3 & 6  &  40   &   0.16 & 39.06 & 0.941  & 0.886\\
4 & 8  &  46   &  31.36 & 45.76 & 0.235  & 0.055\\
5 & 10 &  52   & 134.56 & 52.47 & -0.471 & 0.221\\
\midrule
\hline
$\sum$&   &  & 315.20  &  &  & 9.412\\
\bottomrule
\end{tabular}
\end{center}
Por lo que el $R^2$ es:
$$R^2 = 1-\frac{9.412}{315.20}=0.970$$
y el $R^2_{adj}$ es
$$R^2_{adj} = 1- \frac{4}{3}\frac{9.412}{315.20}=0.96$$
y el error estandar del residuo es:
$$\sigma^2=\frac{SSR}{n-k-1}=\frac{9.412}{3}=3.14$$
Y el error estandar del estimador $\beta_1$ es:
$$s.e.(\hat{\beta_1})=\sqrt{\frac{\sigma^2}{Var(X)}}=\sqrt{\frac{3.14}{27.20/4}}=0.6795$$
    \item ¿Cuándo omitir \(X_2\) sesga \(\hat{\beta}_1\)?
\\
\textbf{Solución}
\\
Cuando $X_2$ tiene relación con $X_1$ y con $Y$.

    \item Estime \(Y_i=\beta_0+\beta_1X_{1i}+\beta_2X_{2i}+u_i\) e informe
      \(\hat{\beta}_0,\hat{\beta}_1,\hat{\beta}_2, R^{2}, R^{2}_{\text{aj}}\) y
      sus errores estándar, sabiendo que:
      \[
(X'X)^{-1}=
\begin{pmatrix}
2.7 & 0.3 & -0.65 \\
0.3 & \;\;0.325 & \;\;-0.35 \\
-0.65 & \;\;.0.35 & \;\;0.425 \\
\end{pmatrix}
\]

\\
\textbf{Solución}
\\
La matriz $X$ esta dada por:

\[
X=\begin{pmatrix}
1 & 4 & 4\\
1 & 4 & 5\\
1 & 6 & 7\\
1 & 8 & 9\\
1 & 10& 9\\
\end{pmatrix},
\qquad
Y=
\begin{pmatrix}
30\\ 34\\ 40 \\ 46 \\ 52
\end{pmatrix}
\]
Entonces

\[
X'Y=\begin{pmatrix}202\\1384\\1452\end{pmatrix},
\qquad
\hat\beta=(X'X)^{-1}X'Y=
\begin{pmatrix}
16.8\\ 2.2\\ 1.4
\end{pmatrix}.
\]
Por lo que $\hat{\beta_0}=16.8$, $\hat{\beta_1}=2.2$ y $\hat{\beta_2}=1.4$. \\
Calculando $R^2$: 
\begin{table}[H]
\centering
\begin{tabular}{|c|c|c|c|c|}
\hline
$i$ &
$Y_i$ (real) &
$\hat Y_i$ (predicho) &
$u_i = Y_i-\hat Y_i$ &
$u_i^{\,2}$ \\ 
\hline
1 & 30 & 31.20 & $-1.20$ & 1.44 \\
2 & 34 & 32.60 & \phantom{$-$}1.40 & 1.96 \\
3 & 40 & 39.80 & \phantom{$-$}0.20 & 0.04 \\
4 & 46 & 47.00 & $-1.00$ & 1.00 \\
5 & 52 & 51.40 & \phantom{$-$}0.60 & 0.36 \\[2pt]
\hline
\multicolumn{4}{r|}{$\displaystyle\sum u_i^{\,2}$} & 4.80 \\
\bottomrule
\end{tabular}
\end{table}
Entonces el $R^2$ es:
$$R^2 = 1-\frac{4.8}{315.2}=0.985$$
y el $R^2_{adj}$ esta dado por:
$$R^2_{adj} = 1- \frac{n-1}{n-k-1}\frac{SSR}{SST}=1-\frac{4}{2}\frac{4.8}{315.2}=0.97$$
para calculo del error estandar de los estimadores, primero debemos encontrar el error estandar del residuo, que esta dado por:
$$\sigma^2=\frac{SSR}{n-k-1}=\frac{4.8}{2}=$$
    \item Compare los coeficientes, \(R^{2}\) y errores estándar de los modelos
      (a) y (c). Interprete.
  \end{enumerate}
\end{enumerate}


\section{Formulario}


\begin{itemize}

\item  Estimadores MCO para RLS.
{
$$\hat{\beta}_0 = \bar{y} - \hat{\beta}_1 \bar{x} $$
$$\hat{\beta}_1 = \frac{\sum_{i=1}^{n}(x_i - \bar{x})(y_i-\bar{y})}{\sum_{i=1}^{n}(x_i - \bar{x})^2}=\frac{Cov(x,y)}{Var(x)}$$}


\item  Estimadores MCO para RLM.
{
$${\hat{\beta}} = (X^TX)^{-1} X^TY $$}


    \item \[ R^2  = 1 - \frac{SSR}{SST} \] 
\item  $$SST = \sum_{i=1}^{n} (y_i - \bar{y})^2 \quad \quad SSR = \sum_{i=1}^{n}  \hat{u}_i^2 $$
    \item \[ \bar{R}^2 =  1 - \frac{(n-1)SSR}{(n-k-1)SST} \]

    \item  $$\hat{\sigma}^2 = \frac{\sum_{i=1}^{n} \hat{u}_i^2 }{n-k-1}= \frac{SSR}{n - k - 1}$$
    \item \[
\text{Var}(\hat{\beta}|X) = \hat{\sigma}^2(X^{\top}X)^{-1}
\]
    \end{itemize}



\end{document}